\documentclass[11pt,a4paper]{article}
\usepackage[margin=2.2cm]{geometry}
\usepackage{amsmath,amssymb}
\usepackage{booktabs}
\usepackage{enumitem}
\usepackage{xcolor}
\usepackage{hyperref}
\usepackage{fancyhdr}
\usepackage{titlesec}
\usepackage{siunitx}

\setlength{\headheight}{14pt}

\pagestyle{fancy}
\fancyhf{}
\rhead{\textit{Solver Derivation}}
\lhead{\textit{Europa Ice Shell Model}}
\cfoot{\thepage}

\titleformat{\section}{\Large\bfseries}{}{0em}{}
\titleformat{\subsection}{\large\bfseries}{}{0em}{}
\titleformat{\subsubsection}{\normalsize\bfseries\itshape}{}{0em}{}

\newcommand{\dd}{\mathrm{d}}
\newcommand{\pp}{\partial}
\newcommand{\Ra}{\mathrm{Ra}}
\newcommand{\Nu}{\mathrm{Nu}}
\newcommand{\half}{\tfrac{1}{2}}

\title{\textbf{Numerical Methods: Full Derivation}\\[0.3em]
\Large Finite Difference Discretisation, Convection\\
Parameterisation, and Bayesian Constraining\\[0.5em]
\large\textit{From Continuous Equations to Matrix Coefficients}}
\author{}
\date{}

\begin{document}
\maketitle
\thispagestyle{fancy}
\tableofcontents
\newpage

%======================================================================
\section{The Heat Equation in Conservative Form}
\label{sec:governing}
%======================================================================

\subsection{Continuous form}

The temperature field $T(z,t)$ in a radially stratified ice shell
satisfies the energy equation with volumetric tidal heating:
\begin{equation}
\rho(T)\,c_p(T)\,\frac{\pp T}{\pp t}
  = \frac{1}{\mathcal{G}(z)}\,\frac{\pp}{\pp z}
    \left[\mathcal{G}(z)\,k(T)\,\frac{\pp T}{\pp z}\right]
  + q_{\mathrm{tidal}}(T)
\label{eq:heat}
\end{equation}
where $z$ is depth from the surface, and the geometric factor accounts
for coordinate system:
\begin{equation}
\mathcal{G}(z) = \begin{cases}
  1       & \text{Cartesian (thin shell, } H < 30\text{~km)} \\
  r(z)^2  & \text{spherical (thick shell, } H \geq 30\text{~km)}
\end{cases}
\end{equation}
with $r(z) = R_p - z$ the radial distance from Europa's centre.

The equation is in \textbf{flux-conservative form}: the divergence
operator $\frac{1}{\mathcal{G}}\frac{\pp}{\pp z}(\mathcal{G}\,k\,\frac{\pp T}{\pp z})$
ensures that the discrete heat flux is conserved across cell interfaces,
even when $k$ jumps by orders of magnitude at the conductive--convective
boundary.

\subsection{Boundary conditions}

\begin{itemize}[leftmargin=2em]
\item \textbf{Surface} ($z = 0$): fixed temperature $T(0,t) = T_{\mathrm{surf}}$
      (Dirichlet).
\item \textbf{Base} ($z = H(t)$): pressure-dependent melting point
      $T(H,t) = T_m(H) = 273.0 - 7.4\times10^{-8}\,\rho\,g\,H$ (Dirichlet).
\end{itemize}

The shell thickness $H(t)$ evolves via the Stefan condition:
\begin{equation}
\frac{\dd H}{\dd t} = \frac{q_{\mathrm{cond}}(H) - q_{\mathrm{ocean}}}{\rho\,L}
\label{eq:stefan}
\end{equation}

\subsection{Initial condition}

The temperature is initialised as a linear profile:
\begin{equation}
T(z, 0) = T_{\mathrm{surf}} + \frac{T_m(H_0) - T_{\mathrm{surf}}}{H_0}\,z
\end{equation}
with initial shell thickness $H_0 = 10$~km.

%======================================================================
\section{Spatial Discretisation}
\label{sec:spatial}
%======================================================================

\subsection{Uniform grid}

The domain $[0, H(t)]$ is divided into $N$ equally spaced nodes
$z_i = i\,\Delta z$ for $i = 0, 1, \ldots, N{-}1$, where:
\begin{equation}
\Delta z = \frac{H(t)}{N - 1}
\end{equation}

The grid uses normalised coordinates $\xi_i = i/(N{-}1) \in [0,1]$
so that physical positions $z_i = \xi_i \cdot H(t)$ automatically
adjust when $H$ changes. This avoids explicit remeshing.

\subsection{Half-node conductivity: harmonic mean}

At cell interfaces, the conductivity $k_{i+\half}$ is computed as the
\textbf{harmonic mean} of adjacent nodal values:
\begin{equation}
k_{i+\half} = \frac{2\,k_i\,k_{i+1}}{k_i + k_{i+1}}
\label{eq:harmonic}
\end{equation}

This choice is physically motivated: for heat flow through two
resistances in series (as across a cell interface), the harmonic mean
gives the correct effective conductivity. It is critical at the
conductive--convective interface, where $k$ may jump by a factor of
$\Nu \sim 10$--$20$ when the Nusselt enhancement is applied. An
arithmetic mean would overestimate the interface flux; the harmonic
mean correctly limits it to the smaller of the two conductivities.

\subsection{Geometric factors at half-nodes}

In spherical geometry, the geometric factor at half-nodes is
approximated by arithmetic averaging:
\begin{equation}
\mathcal{G}_{i+\half} = \frac{\mathcal{G}_i + \mathcal{G}_{i+1}}{2}
  = \frac{r_i^2 + r_{i+1}^2}{2}
\end{equation}

\subsection{Discrete diffusion operator}

Applying the conservative finite-difference stencil to the flux
divergence term at interior node $i$:
\begin{equation}
\frac{1}{\mathcal{G}_i}\frac{\pp}{\pp z}\!\left[\mathcal{G}\,k\,\frac{\pp T}{\pp z}\right]_i
\approx
\frac{1}{\mathcal{G}_i\,\Delta z^2}\left[
  \mathcal{G}_{i+\half}\,k_{i+\half}\,(T_{i+1} - T_i)
  - \mathcal{G}_{i-\half}\,k_{i-\half}\,(T_i - T_{i-1})
\right]
\label{eq:Ldiscrete}
\end{equation}

Define the shorthand coefficients:
\begin{align}
F_i^{+} &= \frac{\mathcal{G}_{i+\half}\,k_{i+\half}}
  {\mathcal{G}_i\,\Delta z^2}
\label{eq:Fplus} \\
F_i^{-} &= \frac{\mathcal{G}_{i-\half}\,k_{i-\half}}
  {\mathcal{G}_i\,\Delta z^2}
\label{eq:Fminus}
\end{align}

Then the discrete spatial operator is:
\begin{equation}
\mathcal{L}_i[T] = F_i^{+}\,(T_{i+1} - T_i) - F_i^{-}\,(T_i - T_{i-1})
  = F_i^{-}\,T_{i-1} - (F_i^{+} + F_i^{-})\,T_i + F_i^{+}\,T_{i+1}
\label{eq:Loperator}
\end{equation}

%======================================================================
\section{Time Integration: the $\theta$-Method}
\label{sec:time}
%======================================================================

\subsection{General $\theta$-method}

The semi-discrete equation
$\rho_i\,c_{p,i}\,\frac{\dd T_i}{\dd t} = \mathcal{L}_i[T] + q_i$
is integrated in time using the $\theta$-weighted scheme:
\begin{equation}
\rho_i\,c_{p,i}\,\frac{T_i^{n+1} - T_i^n}{\Delta t_{\mathrm{eff}}}
  = \theta\,\mathcal{L}_i[T^{n+1}]
  + (1-\theta)\,\mathcal{L}_i[T^n]
  + q_{\mathrm{tidal},i}
\label{eq:theta}
\end{equation}

\begin{itemize}[leftmargin=2em]
\item $\theta = 1$: backward Euler (first-order, L-stable, strongly
      dissipative).
\item $\theta = \half$: Crank--Nicolson (second-order, A-stable, minimally
      dissipative).
\end{itemize}

Define the local CFL-like ratio:
\begin{equation}
\alpha_i = \frac{\Delta t_{\mathrm{eff}}}{\rho_i\,c_{p,i}}
\label{eq:alpha}
\end{equation}

\subsection{Rearrangement into tridiagonal form}

Substituting Eq.~\ref{eq:Loperator} into Eq.~\ref{eq:theta} and
collecting terms on $T^{n+1}$ (left-hand side) and $T^n$ (right-hand
side):

\paragraph{Left-hand side (implicit):}
\begin{align}
\underbrace{-\theta\,\alpha_i\,F_i^{-}}_{a_i}\;T_{i-1}^{n+1}
+ \underbrace{\left[1 + \theta\,\alpha_i\,(F_i^{+} + F_i^{-})\right]}_{b_i}\;T_i^{n+1}
+ \underbrace{-\theta\,\alpha_i\,F_i^{+}}_{c_i}\;T_{i+1}^{n+1}
\label{eq:lhs}
\end{align}

The tridiagonal coefficients are therefore:
\begin{align}
\boxed{a_i = -\theta\,\alpha_i\,F_i^{-}}
  &\quad\text{(sub-diagonal)} \label{eq:ai} \\
\boxed{b_i = 1 + \theta\,\alpha_i\,(F_i^{+} + F_i^{-})}
  &\quad\text{(main diagonal)} \label{eq:bi} \\
\boxed{c_i = -\theta\,\alpha_i\,F_i^{+}}
  &\quad\text{(super-diagonal)} \label{eq:ci}
\end{align}

\paragraph{Right-hand side (explicit + source):}
\begin{equation}
d_i = T_i^n
  + (1-\theta)\,\alpha_i\left[
    F_i^{+,n}(T_{i+1}^n - T_i^n) - F_i^{-,n}(T_i^n - T_{i-1}^n)
  \right]
  + \alpha_i\,q_{\mathrm{tidal},i}
\label{eq:rhs}
\end{equation}

where the superscript $n$ on $F^{\pm,n}$ indicates that the explicit
flux coefficients are evaluated at the old-time material properties.
The tidal heating source term enters the right-hand side directly as:
\begin{equation}
S_i = \frac{\Delta t_{\mathrm{eff}} \cdot q_{\mathrm{tidal}}(T_i)}
  {\rho_i\,c_{p,i}}
\label{eq:source}
\end{equation}

\subsection{Special case: backward Euler ($\theta = 1$)}

When $\theta = 1$, the explicit flux term $(1-\theta)(\ldots)$
vanishes, and the right-hand side simplifies to:
\begin{equation}
d_i = T_i^n + S_i \qquad (\theta = 1)
\end{equation}

This is computationally cheaper (no explicit flux evaluation) and is
used during Rannacher startup.

\subsection{Rannacher startup procedure}

The Crank--Nicolson scheme ($\theta = \half$) is second-order accurate
but can produce spurious oscillations when applied to stiff initial
conditions (e.g., a sharp temperature gradient that develops in early
timesteps). The Rannacher startup (Rannacher 1984) stabilises the
scheme by using backward Euler for the first few steps:

\begin{enumerate}[leftmargin=2em]
\item For the first $N_R = 4$ logical timesteps:
  \begin{enumerate}
  \item Set $\theta = 1$ and $\Delta t_{\mathrm{eff}} = \Delta t / 2$.
  \item Assemble the tridiagonal system and solve for $T^*$ (first half-step).
  \item Using $T^*$ as the new initial condition, assemble and solve again
        for $T^{n+1}$ (second half-step).
  \end{enumerate}
  Each logical step therefore comprises two backward-Euler half-steps,
  advancing by a total of $\Delta t$.

\item For all subsequent steps ($n \geq N_R$):
  \begin{itemize}
  \item Set $\theta = \half$ and $\Delta t_{\mathrm{eff}} = \Delta t$.
  \item Assemble and solve the full Crank--Nicolson system.
  \end{itemize}
\end{enumerate}

The backward-Euler half-steps are unconditionally stable and strongly
dissipate high-frequency oscillations, ensuring a smooth handoff to the
Crank--Nicolson scheme.

%======================================================================
\section{Nonlinear Iteration: Picard Method}
\label{sec:picard}
%======================================================================

The material properties $k(T)$, $\rho(T)$, $c_p(T)$, and
$q_{\mathrm{tidal}}(T)$ all depend on temperature, making the system
nonlinear. At each timestep, the following Picard (fixed-point)
iteration is applied:

\begin{enumerate}[leftmargin=2em]
\item Set the initial guess $T^{(0)} = T^n$ (previous timestep solution).
\item For $m = 0, 1, 2$ (maximum 3 iterations):
  \begin{enumerate}
  \item Evaluate material properties at $T^{(m)}$:
    $k^{(m)}_i = k(T^{(m)}_i)$,
    $\rho^{(m)}_i = \rho(T^{(m)}_i)$,
    $c_{p,i}^{(m)} = c_p(T^{(m)}_i)$,
    $q^{(m)}_i = q_{\mathrm{tidal}}(T^{(m)}_i)$.
  \item Compute half-node conductivities $k^{(m)}_{i+\half}$ via the
        harmonic mean (Eq.~\ref{eq:harmonic}).
  \item Assemble the tridiagonal system using coefficients
        Eqs.~\ref{eq:ai}--\ref{eq:ci} and right-hand side Eq.~\ref{eq:rhs},
        with implicit coefficients from $T^{(m)}$ and explicit
        fluxes from $T^n$.
  \item Solve the tridiagonal system via the Thomas algorithm to obtain
        $T^{(m+1)}$.
  \item Check convergence: if
        $\max_i |T^{(m+1)}_i - T^{(m)}_i| < \SI{0.01}{\kelvin}$,
        accept $T^{n+1} = T^{(m+1)}$ and exit.
  \end{enumerate}
\item If convergence is not achieved after 3 iterations, accept the
      last iterate as $T^{n+1}$.
\end{enumerate}

In practice, convergence is typically achieved in 1--2 iterations
because the temperature change per timestep is small relative to the
absolute temperature.

%======================================================================
\section{Tridiagonal Solution: Thomas Algorithm}
\label{sec:thomas}
%======================================================================

The assembled system $\mathbf{A}\,\mathbf{T}^{n+1} = \mathbf{d}$ is a
tridiagonal matrix equation of dimension $N$:
\begin{equation}
\begin{pmatrix}
b_0 & c_0 \\
a_1 & b_1 & c_1 \\
    & a_2 & b_2 & c_2 \\
    &     & \ddots & \ddots & \ddots \\
    &     &        & a_{N-2} & b_{N-2} & c_{N-2} \\
    &     &        &         & a_{N-1} & b_{N-1}
\end{pmatrix}
\begin{pmatrix} T_0 \\ T_1 \\ T_2 \\ \vdots \\ T_{N-2} \\ T_{N-1} \end{pmatrix}
=
\begin{pmatrix} d_0 \\ d_1 \\ d_2 \\ \vdots \\ d_{N-2} \\ d_{N-1} \end{pmatrix}
\end{equation}

This is solved via \texttt{scipy.linalg.solve\_banded} (the Thomas
algorithm), which performs LU decomposition in $O(N)$ operations:

\paragraph{Forward sweep} (eliminate sub-diagonal):
\begin{align}
w_i &= \frac{a_i}{b_{i-1}'} \\
b_i' &= b_i - w_i\,c_{i-1} \\
d_i' &= d_i - w_i\,d_{i-1}'
\end{align}

\paragraph{Back substitution:}
\begin{align}
T_{N-1} &= \frac{d_{N-1}'}{b_{N-1}'} \\
T_i     &= \frac{d_i' - c_i\,T_{i+1}}{b_i'} \qquad i = N{-}2, \ldots, 0
\end{align}

%======================================================================
\section{Boundary Condition Implementation}
\label{sec:bc_impl}
%======================================================================

\subsection{Surface node ($i = 0$): Dirichlet}

For a fixed surface temperature, the first row of the tridiagonal
system is:
\begin{equation}
\underbrace{1}_{b_0}\,T_0^{n+1} + \underbrace{0}_{c_0}\,T_1^{n+1}
= \underbrace{T_{\mathrm{surf}}}_{d_0}
\end{equation}

\subsection{Surface node ($i = 0$): Stefan--Boltzmann radiative}

For the radiative boundary condition
$\varepsilon\sigma T_s^4 = (1-a)Q + k\,\frac{T_1 - T_s}{\Delta z}$,
the $T^4$ nonlinearity is linearised around the current temperature
estimate $T_s^{(m)}$:
\begin{equation}
T^4 \approx 4\,(T_s^{(m)})^3\,T_s^{n+1} - 3\,(T_s^{(m)})^4
\end{equation}

Substituting and rearranging into the matrix row
$b_0\,T_0^{n+1} + c_0\,T_1^{n+1} = d_0$:
\begin{align}
b_0 &= 4\,\varepsilon\sigma\,(T_s^{(m)})^3 + \frac{k_{\half}}{\Delta z}
  \label{eq:b0_rad} \\
c_0 &= -\frac{k_{\half}}{\Delta z} \label{eq:c0_rad} \\
d_0 &= (1-a)\,Q_{\odot} + 3\,\varepsilon\sigma\,(T_s^{(m)})^4
  \label{eq:d0_rad}
\end{align}

where $k_{\half}$ is the half-node conductivity between nodes 0 and 1.

\subsection{Basal node ($i = N{-}1$): melting-point Dirichlet}

The last row enforces the pressure-dependent melting temperature:
\begin{equation}
\underbrace{1}_{b_{N-1}}\,T_{N-1}^{n+1} = T_m(H)
  = 273.0 - 7.4\times10^{-8}\,\rho\,g\,H
\end{equation}

\subsection{Stefan condition: shell thickness evolution}

After solving for $T^{n+1}$, the basal heat flux is computed via a
second-order one-sided difference:
\begin{equation}
q_{\mathrm{cond}} = k(T_{N-1})\,
  \frac{3\,T_{N-1} - 4\,T_{N-2} + T_{N-3}}{2\,\Delta z}
\label{eq:qcond_numerical}
\end{equation}

The freezing-front velocity is:
\begin{equation}
v = \frac{q_{\mathrm{cond}} - q_{\mathrm{ocean}}}{\rho\,L}
\end{equation}

The shell thickness is updated explicitly:
\begin{equation}
H^{n+1} = \max\!\left(H^n + v\,\Delta t_{\mathrm{eff}},\; 500\text{~m}\right)
\end{equation}

The 500~m floor prevents the shell from vanishing entirely, which
would produce a singular grid. Equilibrium is declared when
$|v| < 10^{-15}$~m/s.

Because the grid is defined on normalised coordinates
$\xi_i \in [0,1]$, the physical node positions $z_i = \xi_i\,H$
and spacing $\Delta z = H/(N{-}1)$ automatically adjust when $H$
changes. No explicit remeshing or interpolation is required.

%======================================================================
\section{Convection Parameterisation: Full Algorithm}
\label{sec:convection}
%======================================================================

\subsection{Physical basis}

Europa's ice shell is expected to convect in the stagnant-lid regime:
a cold, rigid conductive lid overlies a slowly overturning convective
interior. The model does not resolve convective flow explicitly.
Instead, it follows the parameterised approach of Green et al.\ (2021)
and Deschamps \& Vilella (2021), which replaces the convective heat
transport with an enhanced effective conductivity.

The algorithm has four stages:
\begin{enumerate}[leftmargin=2em]
\item Compute the rheological transition temperature $T_c$.
\item Scan the temperature profile to find the interface depth.
\item Compute Ra and Nu for the convective sublayer.
\item Modify the conductivity profile.
\end{enumerate}

\subsection{Stage 1: Transition temperature $T_c$}

The transition temperature marks where viscosity becomes low enough
for convective participation.

\subsubsection{Interior temperature (Deschamps \& Vilella 2021, Eq.~18)}

The mean temperature of the convective interior is:
\begin{equation}
T_i = B\left[\sqrt{1 + \frac{2}{B}\left(T_m - c_2\,\Delta T_s\right)} - 1\right]
\label{eq:Ti}
\end{equation}
where:
\begin{equation}
B = \frac{Q_v}{2\,R\,c_1}, \qquad
\Delta T_s = T_m - T_{\mathrm{surf}}
\end{equation}
with $c_1 = 1.43$ and $c_2 = -0.03$ (Deschamps \& Vilella 2021),
$Q_v = 59{,}400$~J/mol, and $R = 8.314$~J/(mol$\cdot$K).

\subsubsection{Viscous temperature scale}

The Frank--Kamenetskii parameter characterises the exponential viscosity
contrast:
\begin{equation}
A = \frac{Q_v}{N\,R\,T_m}
\end{equation}
where $N = 1$ for Newtonian (diffusion) creep. The viscosity at $T_i$
and its derivative are:
\begin{align}
\eta(T_i) &= \eta_{\mathrm{ref}}\,\exp\!\left[A\left(\frac{T_m}{T_i} - 1\right)\right]
\label{eq:eta_Ti} \\
\frac{\dd\eta}{\dd T}\bigg|_{T_i} &= -\eta(T_i)\,\frac{A\,T_m}{T_i^2}
\label{eq:deta}
\end{align}

The viscous temperature scale is:
\begin{equation}
\Delta T_v = -\frac{\eta(T_i)}{\dd\eta/\dd T\big|_{T_i}}
  = \frac{T_i^2}{A\,T_m}
  = \frac{N\,R\,T_i^2}{Q_v}
\label{eq:DTv}
\end{equation}

\subsubsection{Lid-base temperature}

The conductive lid extends from the surface down to the depth where
the temperature contrast from the convective interior equals
$\theta_{\mathrm{lid}}$ viscous temperature scales:
\begin{equation}
\boxed{T_c = T_i - \theta_{\mathrm{lid}}\,\Delta T_v}
\label{eq:Tc}
\end{equation}
with $\theta_{\mathrm{lid}} = 2.24$ (Green et al.\ 2021).

\subsubsection{Alternative: Howell (2021) method}

A simpler analytical estimate sets:
\begin{equation}
T_c^{\mathrm{Howell}}
  = \frac{\sqrt{4\,T_m\,(R/Q_v) + 1} - 1}{2\,(R/Q_v)}
\end{equation}
and the conductive base temperature as $T_{\mathrm{cond,base}} = 2\,T_c - T_m$.
This is faster but less accurate than the Green method for strongly
internally heated shells.

\subsection{Stage 2: Interface depth from temperature profile}

Given $T_c$, the algorithm scans the numerical temperature profile
$T_0, T_1, \ldots, T_{N-1}$ from the surface:

\begin{enumerate}[leftmargin=2em]
\item Find the first index $j$ such that $T_j \geq T_c$.
\item If $0 < j < N{-}1$, linearly interpolate to find the exact
      interface depth:
\begin{equation}
z_c = z_{j-1} + \frac{T_c - T_{j-1}}{T_j - T_{j-1}}\,(z_j - z_{j-1})
\label{eq:zc_interp}
\end{equation}
\item Set $D_{\mathrm{cond}} = z_c$ and $D_{\mathrm{conv}} = H - z_c$.
\end{enumerate}

If no node satisfies $T_j \geq T_c$, the entire shell is conductive
($D_{\mathrm{conv}} = 0$, $\Nu = 1$).

\subsection{Stage 3: Rayleigh and Nusselt numbers}

\subsubsection{Rayleigh number}

The Rayleigh number of the convective sublayer is:
\begin{equation}
\Ra = \frac{\rho\,g\,\alpha_{th}\,\Delta T_{\mathrm{conv}}\,
  D_{\mathrm{conv}}^3}{\kappa\,\bar{\eta}}
\label{eq:Ra}
\end{equation}
where:
\begin{align}
\Delta T_{\mathrm{conv}} &= T_m - T_c \\
\bar{T} &= \tfrac{1}{2}(T_m + T_c) \\
\rho &= \rho(\bar{T}) = 917\,[1 + 1.6\times10^{-4}\,(273 - \bar{T})] \\
k &= k(\bar{T}) = 567/\bar{T} \\
c_p &= c_p(\bar{T}) = 7.49\,\bar{T} + 90 \\
\kappa &= k/(\rho\,c_p) \quad\text{(thermal diffusivity)} \\
\bar{\eta} &= \eta(\bar{T}, d_{\mathrm{grain}})
  \quad\text{(composite viscosity, Eq.~\ref{eq:viscosity_repeat})}
\end{align}

Convection is active if and only if $\Ra \geq \Ra_{\mathrm{crit}} = 1000$.

\subsubsection{Nusselt number (Green et al.\ 2021)}

For $\Ra \geq \Ra_{\mathrm{crit}}$:
\begin{equation}
\boxed{\Nu = C\,\Ra^{\xi}\left(\frac{T_i - T_c}
  {\Delta T_{\mathrm{conv}}}\right)^{\!\zeta}}
\label{eq:Nu_green}
\end{equation}
with $C = 0.3446$, $\xi = 1/3$ (Solomatov \& Moresi 2000), and
$\zeta = 4/3$ (internal heating correction; Green et al.\ 2021).

The ratio $(T_i - T_c)/\Delta T_{\mathrm{conv}}$ accounts for the
fact that internal tidal heating raises the interior temperature $T_i$
above what would be expected from purely basal heating. When
$T_i > T_c$ (always true by construction), this ratio is $> 0$ and
amplifies the Nusselt number. It is clipped to a minimum of 0.01
for numerical safety.

For $\Ra < \Ra_{\mathrm{crit}}$: $\Nu = 1$ (pure conduction).

\subsection{Stage 4: Conductivity profile modification}

The conductivity profile used in the tridiagonal assembly is modified
as follows:

\begin{enumerate}[leftmargin=2em]
\item Start with the base conductivity at each node:
\begin{equation}
k_i^{\mathrm{base}} = \frac{567}{T_i}
\end{equation}

\item Apply porosity reduction in the cold upper shell:
\begin{equation}
k_i = \begin{cases}
  k_i^{\mathrm{base}}\,(1 - f_{\mathrm{porosity}}) & T_i < T_\phi \\
  k_i^{\mathrm{base}} & T_i \geq T_\phi
\end{cases}
\end{equation}

\item Apply the Nusselt enhancement below the conductive--convective
      interface. For all nodes at or below the interface index $j_c$:
\begin{equation}
k_i \leftarrow \Nu_{\mathrm{eff}} \cdot k_i
  \qquad \text{for } i \geq j_c
\label{eq:kmod}
\end{equation}
where:
\begin{equation}
\Nu_{\mathrm{eff}} = \Nu
\label{eq:Nu_eff}
\end{equation}

\item Compute half-node conductivities via harmonic mean
      (Eq.~\ref{eq:harmonic}) for the tridiagonal assembly.
\end{enumerate}

\subsection{Shared convection closure}

The current 2D implementation uses the same stagnant-lid Nusselt
enhancement as the 1D solver in every latitude column, with no extra
surface-temperature ramp. Convective suppression therefore occurs only
through the temperature-profile scan and the Rayleigh criticality test,
not through an additional latitude-specific closure.

%======================================================================
\section{Complete Timestep Algorithm}
\label{sec:algorithm}
%======================================================================

Each timestep proceeds as follows:

\begin{enumerate}[leftmargin=2em]
\item \textbf{Determine time integration parameters.}
  If $n < N_R = 4$: set $\theta = 1$, $\Delta t_{\mathrm{eff}} = \Delta t/2$.
  Otherwise: set $\theta = \half$, $\Delta t_{\mathrm{eff}} = \Delta t$.

\item \textbf{Begin Picard loop} ($m = 0, 1, 2$):
  \begin{enumerate}
  \item Evaluate $k^{(m)}_i$, $\rho^{(m)}_i$, $c_{p,i}^{(m)}$,
        $q_{\mathrm{tidal},i}^{(m)}$ at $T^{(m)}$.
  \item \textbf{Build conductivity profile:}
    \begin{itemize}
    \item Compute $T_c$ via Green/Deschamps (Eqs.~\ref{eq:Ti}--\ref{eq:Tc}).
    \item Scan profile for interface depth $z_c$ (Eq.~\ref{eq:zc_interp}).
    \item Compute Ra (Eq.~\ref{eq:Ra}) and Nu (Eq.~\ref{eq:Nu_green}).
    \item Apply porosity and Nusselt modifications to $k$ profile.
    \item Compute harmonic-mean half-node conductivities.
    \end{itemize}
  \item \textbf{Assemble tridiagonal system:}
    compute $a_i$, $b_i$, $c_i$ (Eqs.~\ref{eq:ai}--\ref{eq:ci})
    and $d_i$ (Eq.~\ref{eq:rhs}).
  \item \textbf{Apply boundary conditions:} set row 0 (surface) and
        row $N{-}1$ (base).
  \item \textbf{Solve} via Thomas algorithm $\to T^{(m+1)}$.
  \item \textbf{Check convergence:}
        $\max_i |T_i^{(m+1)} - T_i^{(m)}| < 0.01$~K.
  \end{enumerate}

\item \textbf{If Rannacher step} ($n < N_R$): repeat step 2 for the
      second half-step using $T^{(m+1)}$ as the new initial condition.

\item \textbf{Update shell thickness} via Stefan condition
      (Eq.~\ref{eq:stefan}).

\item \textbf{Check equilibrium:} if $|\dd H/\dd t| < 10^{-15}$~m/s, halt.

\item Increment $n$ and return to step 1.
\end{enumerate}

%======================================================================
\section{Bayesian Importance Reweighting}
\label{sec:bayes}
%======================================================================

\subsection{Motivation}

The forward model maps sampled parameters $\boldsymbol{\theta}_i$ to an
equilibrium conductive lid thickness $D_{\mathrm{cond},i}$. After
running $N$ Monte Carlo realisations, we wish to update the ensemble
to reflect the Juno MWR observation
$D_{\mathrm{cond}}^{\mathrm{obs}} = 29 \pm 10$~km without re-running
the expensive forward model.

\subsection{Likelihood}

The likelihood of realisation $i$ given the observation is:
\begin{equation}
\mathcal{L}_i = P(D^{\mathrm{obs}} \mid D_{\mathrm{cond},i})
  = \frac{1}{\sqrt{2\pi}\,\sigma_{\mathrm{tot}}}\,
  \exp\!\left[-\frac{1}{2}
  \left(\frac{D_{\mathrm{cond},i} - D^{\mathrm{obs}}}
    {\sigma_{\mathrm{tot}}}\right)^{\!2}\right]
\end{equation}
where:
\begin{equation}
\sigma_{\mathrm{tot}} = \sqrt{\sigma_{\mathrm{obs}}^2 + \sigma_{\mathrm{model}}^2}
  = \sqrt{10^2 + 3^2} \approx 10.4\text{~km}
\end{equation}
The model discrepancy $\sigma_{\mathrm{model}} = 3$~km accounts for
numerical resolution and parameterisation uncertainty.

\subsection{Importance weights}

Each realisation receives a normalised importance weight:
\begin{equation}
w_i = \frac{\mathcal{L}_i}{\sum_{j=1}^N \mathcal{L}_j}
\end{equation}

In practice, to avoid numerical underflow, the log-weights are computed
first:
\begin{equation}
\ln w_i = -\frac{1}{2}\left(\frac{D_{\mathrm{cond},i} - D^{\mathrm{obs}}}
  {\sigma_{\mathrm{tot}}}\right)^{\!2}
\end{equation}
then stabilised by subtracting the maximum before exponentiating:
\begin{equation}
w_i = \frac{\exp(\ln w_i - \max_j \ln w_j)}
  {\sum_k \exp(\ln w_k - \max_j \ln w_j)}
\end{equation}

\subsection{Effective sample size}

The quality of the reweighting is monitored via the effective sample
size:
\begin{equation}
\mathrm{ESS} = \frac{\left(\sum_i w_i\right)^2}{\sum_i w_i^2}
  = \frac{1}{\sum_i \tilde{w}_i^2}
\end{equation}
where $\tilde{w}_i = w_i / \sum_j w_j$ are the normalised weights.

An ESS fraction $\mathrm{ESS}/N > 10\%$ indicates acceptable
prior--likelihood overlap. If $\mathrm{ESS}/N < 10\%$, the prior
is too diffuse relative to the likelihood, and the posterior is
dominated by a few high-weight samples.

\subsection{Weighted posterior statistics}

Posterior percentiles are computed from the weighted empirical CDF.
Sort the $D_{\mathrm{cond},i}$ values and their weights, compute the
cumulative weight, and interpolate:
\begin{equation}
\hat{q}_p: \quad \sum_{j: D_j \leq \hat{q}_p} w_j = p
\end{equation}

The posterior median ($p = 0.5$), 68\% interval ($p = 0.16, 0.84$),
and 90\% interval ($p = 0.05, 0.95$) are reported.

\subsection{Iterative prior tightening}

When the initial prior is broad, the posterior may be poorly resolved.
The model implements an iterative scheme:

\begin{enumerate}[leftmargin=2em]
\item Run Monte Carlo with unconstrained (audited) priors.
\item Compute importance weights and posterior statistics.
\item Derive tightened priors from the posterior:
  \begin{itemize}
  \item $q_{\mathrm{basal}}$: new uniform bounds from posterior
        [5th, 95th] percentiles.
  \item $d_{\mathrm{grain}}$: new lognormal centre from posterior median,
        new $\sigma$ from posterior 68\% log-width.
  \end{itemize}
\item Re-run Monte Carlo with tightened priors.
\item Repeat until the posterior median falls within the target range
      and $\mathrm{ESS}/N > 10\%$.
\end{enumerate}

\subsection{Bayes factors for model comparison}

To compare ocean transport scenarios $\mathcal{M}_j$ and
$\mathcal{M}_k$, the Bayes factor is:
\begin{equation}
\mathrm{BF}_{jk}
  = \frac{P(D^{\mathrm{obs}} \mid \mathcal{M}_j)}
         {P(D^{\mathrm{obs}} \mid \mathcal{M}_k)}
  = \frac{\frac{1}{N_j}\sum_i \mathcal{L}_i^{(j)}}
         {\frac{1}{N_k}\sum_i \mathcal{L}_i^{(k)}}
\end{equation}

The marginal likelihood for each scenario is the mean likelihood
across its ensemble, estimated by the arithmetic mean of the
(unnormalised) weights.

%======================================================================
\section{Composite Viscosity Reference}
\label{sec:visc_ref}
%======================================================================

For completeness, the composite diffusion-creep viscosity used in
Ra computation and tidal heating:
\begin{equation}
\eta(T, d) = \frac{1}{2}\left[
  \frac{42\,\Omega}{R\,T\,d^2}
  \left(
    D_{0v}\,e^{-Q_v/(RT)}
    + \frac{\pi\,\delta}{d}\,D_{0b}\,e^{-Q_b/(RT)}
  \right)
\right]^{-1}
\label{eq:viscosity_repeat}
\end{equation}

clipped to $[10^{12},\, 10^{25}]$~Pa$\cdot$s.

\end{document}
